% Book pp. 111-116 (PDF pp. 112-117)
\section{Applying Descartes' Rule of Signs Theorem}

There is another way that can be used to determine how many positive and negative
real zeros a polynomial function might have. To do this, first write the polynomial
in descending order (from the highest degree to the lowest) if it is not already. We
use \textbf{Descartes' Rule of Signs}. This rule helps you see how the number of times
the signs change in the polynomial is related to the number of positive real zeros.

For example, if you look at the polynomial function below, you will notice it has
one sign change, which gives you information about its positive zeros.

Let us go through an example using Descartes' Rule of Signs to determine the
possible number of positive and negative real zeros.

\begin{example}{Example 3.9}
Consider the polynomial:

$f(x) = 3x^3 - 2x^2 + 4x - 5$

\solution
\textbf{Step 1:} \textit{Identify Sign Changes for Positive Real Zeros}
\begin{enumerate}
  \item Write the polynomial in descending order:
  \item $f(x) = 3x^3 - 2x^2 + 4x + 5$
  \booknote[S3-10]{the polynomial is $3x^3 - 2x^2 + 4x - 5$, as in the question
  and in the coefficient list that follows; the ``$+ 5$'' is a misprint.}

  Look at the coefficients:

  3 (positive)

  --2 (negative)

  4 (positive)

  $-5$ (negative)
  \item Determine the sign changes:

  From 3 to -2 (1 change)

  From $-2$ to 4 (2 changes)

  From 4 to $-5$ (3 changes)

  Total sign changes for positive zeros $=$ 3. This means there could be 3,
  (3 -- 2 $=$ 1) positive real zeros. Note, the number of possible roots is found by
  subtracting 2 each time until we would get a negative number.
\end{enumerate}
\textbf{Step 2:} \textit{Identify Sign Changes for Negative Real Zeros}
\begin{quote}
Now, let's find the negative real zeros by evaluating f($-x$):

$f(-x) = 3(-x)^3 - 2(-x)^2 + 4(-x) - 5$

$f(-x) = -3x^3 - 2x^2 - 4x - 5$
\end{quote}
\textbf{Step 3:} \textit{Look at the Signs of f($-x$)}
\begin{quote}
The coefficients are:

$-3$ (negative)

$-2$ (negative)

$-4$ (negative)

$-5$ (negative)
\end{quote}
\textbf{Step 4:} \textit{Determine Sign Changes for Negative Real Zeros}
\begin{quote}
There are \textbf{no sign changes} in f($-$x)
\end{quote}
Total sign changes for negative zeros $=$ 0. This means there are no negative real
roots or zeros.
\end{example}

This example illustrates how to use Descartes' Rule of Signs to analyse a polynomial
and determine the possible numbers of positive and negative real zeros.

Let us solve more examples.

\begin{example}{Example 3.10}
Use Descartes' Rule of Signs to determine how many possible positive and
negative real zeros of the following:
\begin{enumerate}
  \item[a.] $g(x) = 2x^4 - 11x^3 + 4x^2 - 8x - 40$
  \item[b.] $g(x) = 9x^3 - 10x^2 + 5x + 3$
  \item[c.] f$(x) = -x^5 + 3x^4 - 5x^3 - 12x^2 + 6x + 4$
\end{enumerate}

\solution
\begin{enumerate}
  \item[a.] $g(x) = 2x^4 - 11x^3 + 4x^2 - 8x - 40$
\end{enumerate}
\textbf{Step 1:} \textit{Identify Sign Changes for Positive Real Zeros}
\begin{quote}
Write the polynomial in descending order:

$g(x) = 2x^4 - 11x^3 + 4x^2 - 8x - 40$

Look at the coefficients:

2 (positive)

--11 (negative)

4 (positive)

$-8$ (negative)

$-40$ (negative)

Look at the coefficients:

Determine the sign changes:

From 2 to -11 (1 change)

From $-11$ to 4 (2 changes)

From 4 to $-8$ (3 changes)

From -8 to $-40$ (no change)

Total sign changes for positive zeros $=$ 3. This means there could be 3 or 1
positive real zeros.
\end{quote}
\textbf{Step 2:} \textit{Identify Sign Changes for Negative Real Zeros}
\begin{quote}
Now, let's find the negative real zeros by evaluating f($-$x):

$g(x) = 2(-x)^4 - 11(-x)^3 + 4(-x)^2 - 8(-x) - 40$

$g(x) = 2x^4 + 11x^3 + 4x^2 + 8x - 40$
\end{quote}
\textbf{Step 3:} \textit{Look at the Signs of f($-x$)}
\begin{quote}
The coefficients are:

2 (positive)

11 (positive)

4 (positive)

8 (positive)

$-40$ (negative)
\end{quote}
\textbf{Step 4:} \textit{Determine Sign Changes for Negative Real Zeros}
\begin{quote}
Determine the sign changes:

From 2 to 11 (no change)

From 11 to 4 (no changes)

From 4 to 8 (no changes)

From 8 to $-40$ (1 change)

Total sign changes for negative zeros $=$ 1. This means there is 1 negative
real zero.
\end{quote}
\begin{enumerate}
  \item[b.] $g(x) = 9x^3 - 10x^2 + 5x + 3$
\end{enumerate}
\textbf{Step 1:} \textit{Identify Sign Changes for Positive Real Zeros}
\begin{quote}
Write the polynomial in descending order: $g(x) = 9x^3 - 10x^2 + 5x + 3$.

Look at the coefficients:

9 (positive)

$-10$ (negative)

5 (positive)

3 (negative)
\booknote[S3-11]{3 is positive; the sign changes counted below (2) are right.}

Determine the sign changes:

From 9 to -10 (1 change)

From $-10$ to 5 (2 changes)

From 5 to 3 (no change)

Total sign changes for positive zeros $=$ 2. This means there could be 2 or 0.
\end{quote}
\textbf{Step 2:} \textit{Identify Sign Changes for Negative Real Zeros}
\begin{quote}
Now, let's find the negative real zeros by evaluating g($-$x):

$g(x) = 9(-x)^3 - 10(-x)^2 + 5(-x) + 3$

$g(x) = -9x^3 - 10x^2 - 5x + 3$
\end{quote}
\textbf{Step 3:} \textit{Look at the Signs of f($-x$)}
\begin{quote}
The coefficients are:

$-9$ (negative)

$-10$ (negative)

$-5$ (negative)

3 (positive)
\end{quote}
\textbf{Step 4:} \textit{Determine Sign Changes for Negative Real Zeros}
\begin{quote}
Determine the sign changes:

From -9 to -10 (no change)

From -10 to -5 (no changes)

From -5 to 3 (1 change)

Total sign changes for negative zeros $=$ 1. This means there is 1 negative
real zero.
\end{quote}
\begin{enumerate}
  \item[c.] f$(x) = -x^5 + 3x^4 - 5x^3 - 12x^2 + 6x + 4$
\end{enumerate}
\textbf{Step 1:} \textit{Identify Sign Changes for Positive Real Zeros}
\begin{quote}
Write the polynomial in descending order:

f$(x) = -x^5 + 3x^4 - 5x^3 - 12x^2 + 6x + 4$

Look at the coefficients:

$-1$ (negative)

3 (positive)

$-5$ (negative)

$-12$ (negative)

6 (positive)

4 (positive)

\textit{Determine the sign changes:}

From $-1$ to 3 (1 change)

From 3 to $-5$ (2 changes)

\textbf{From $-5$ to $-12$ (no changes)}

From $-12$ to 6 (3 changes)

From 6 to 4 (no change)

Total sign changes for positive zeros $=$ 3. This means there could be 3 or 1
positive real zeros.
\end{quote}
\textbf{Step 2:} \textit{Identify Sign Changes for Negative Real Zeros}
\begin{quote}
Now, let's find the negative real zeros by evaluating f($-$x):

f$(x) = -(-x)^5 + 3(-x)^4 - 5(-x)^3 - 12(-x)^2 + 6(-x) + 4$

f$(x) = x^5 + 3x^4 + 5x^3 - 12x^2 - 6x + 4$
\end{quote}
\textbf{Step 3:} Look at the Signs of f($-$x)
\begin{quote}
The coefficients are:

1 (positive)

3 (positive)

5 (positive)

--12 (negative)

--6 (negative)

4 (positive)
\end{quote}
\textbf{Step 4:} \textit{Determine Sign Changes for Negative Real Zeros}
\begin{quote}
Determine the sign changes:

From 1 to 3 (no change)

From 3 to 5 (no changes)

From 5 to $-12$ (2 changes)

From $-2$ to $-6$ no changes)
\booknote[S3-12]{``From $-2$'' should read ``From $-12$''; and the count of changes
in the line above (2) should be 1, the line below 2, giving the printed total of 2.}

From $-6$ to 4 (3 change)

Total sign changes for negative zeros $=$ 2. This means there is 2 or 0 negative
real zeros.
\end{quote}
\end{example}
